\documentclass[9pt, spanish, twoside]{article}

% --- Motores de Estilo "Vogue High-Tech" ---
\usepackage[utf8]{inputenc}
\usepackage[T1]{fontenc}
\usepackage[spanish, es-tabla, es-nodecimaldot]{babel}
\usepackage[letterpaper, margin=1cm, top=1.3cm, bottom=1.5cm, headsep=0.7cm]{geometry}
\usepackage{multicol}
\usepackage{xcolor}
\usepackage{titlesec}
\usepackage{fancyhdr}
\usepackage{lettrine}
\usepackage{tcolorbox}
\usepackage{booktabs}
\usepackage{microtype}
\usepackage{enumitem}
\usepackage{amsmath, amssymb}
\usepackage{pgfplots}
\usepackage{animate}
\usepackage{hyperref}
\pgfplotsset{compat=1.18}
\usetikzlibrary{arrows.meta, positioning}

% --- Tipografía de Élite 2026 ---
\usepackage{mathpazo} % Palatino (Serif - NYT Style)
\usepackage[scaled=0.85]{helvet} % Helvetica (Sans - MIT Style)
\newcommand{\techfont}{\fontfamily{phv}\selectfont}

% --- Paleta "Titanium & Champagne" 2026 ---
\definecolor{VogueBlack}{RGB}{8, 8, 8}
\definecolor{NASAClassic}{RGB}{0, 32, 91}
\definecolor{Champagne}{RGB}{181, 158, 118}
\definecolor{TitaniumGray}{RGB}{242, 242, 242}
\definecolor{InkGray}{RGB}{60, 60, 60}
\definecolor{CriticalRed}{RGB}{180, 30, 30}

% --- Configuración de Títulos de Alta Gama ---
\titleformat{\section}
  {\techfont\fontsize{10}{11}\selectfont\bfseries\color{VogueBlack}\lsstyle}
  {}{0em}{\uppercase}
  [\vspace{1pt}\color{Champagne}\hrule height 0.4pt]

\titleformat{\subsection}
  {\techfont\fontsize{7}{8}\selectfont\bfseries\color{InkGray}\lsstyle}
  {}{0em}{\uppercase}

% --- Encabezado "The 2026 Intelligence" ---
\pagestyle{fancy}
\fancyhf{}
\lhead{\scriptsize \techfont \textbf{UPTC ECONOMICS} \textbar \ \textcolor{Champagne}{QUANTITATIVE RESEARCH}}
\rhead{\scriptsize \techfont \textcolor{VogueBlack}{ECONOMETRICS PROTOCOL \textbar \ ENE. 2026}}
\fancyfoot[C]{\scriptsize \techfont \thepage}
\renewcommand{\headrulewidth}{0.3pt}

% --- Comandos de Diseño Pro ---
\newcommand{\vogueQuote}[1]{
    \begin{center}
        \vspace{6pt}
        {\color{Champagne}\hrule height 0.2pt}
        \vspace{5pt}
        {\fontsize{11}{14} \selectfont \itshape \color{VogueBlack} ``#1''}
        \vspace{5pt}
        {\color{Champagne}\hrule height 0.2pt}
        \vspace{6pt}
    \end{center}
}

\newcommand{\critical}[1]{\textcolor{CriticalRed}{\textbf{#1}}}
\newcommand{\eqbox}[1]{\colorbox{TitaniumGray}{$\displaystyle #1$}}

% --- Configuración de Hyperref ---
\hypersetup{
    colorlinks=true,
    linkcolor=NASAClassic,
    urlcolor=Champagne,
    citecolor=InkGray
}

\setlength{\parindent}{0pt}
\setlength{\columnsep}{16pt}

\begin{document}

% --- PORTADA CINEMATOGRÁFICA ---
\thispagestyle{empty}
\begin{center}
    \vspace*{-0.8cm}
    {\techfont \fontsize{6}{7.5} \selectfont \color{InkGray} \lsstyle TUNJA, COLOMBIA \textbar \ UPTC FACULTAD DE CIENCIAS ECONÓMICAS \textbar \ RELEASE: 05.01.2026} \\
    \vspace{10pt}
    \hrule height 0.2pt
    \vspace{30pt}
    {\fontsize{52}{48} \selectfont \textbf{ECONOMETRÍA}} \\
    \vspace{4pt}
    {\techfont \fontsize{8}{10} \selectfont \color{Champagne} \lsstyle CHEATSHEET ULTRA-RIGUROSO DE ALTA GAMA} \\
    \vspace{18pt}
    {\color{Champagne} \hrule height 1pt width 1.5cm}
    \vspace{18pt}
    {\Large \itshape \color{InkGray} Un compendio sistémico de la teoría econométrica moderna, desde los fundamentos clásicos hasta las fronteras asintóticas.} \\
    \vspace{25pt}
    {\techfont \fontsize{7}{9} \selectfont \lsstyle COMPILADO POR \textbf{\uppercase{Emanuel Quintana Silva}} \textbar \ 2026} \\
    \vspace{15pt}
    \hrule height 0.2pt
    \vspace{10pt}
    {\scriptsize \techfont \color{InkGray}
    \textbf{REFERENCIAS CORE:} Gujarati \& Porter (2009) • Wooldridge (2015) • Davidson \& MacKinnon (2004) • Hamilton (1994) • Greene (2018)}
\end{center}

\newpage

% === INICIO DEL CONTENIDO ===
\begin{multicols*}{3}
\setlength{\columnseprule}{0pt}

\section{Definición y Núcleo}
\subsection{Concepto Central}
\lettrine[lines=2, nindent=0pt]{\textbf{\textcolor{Champagne}{L}}}{a} \textbf{econometría} es la ciencia social que aplica teoría económica, matemáticas e inferencia estadística al \critical{análisis cuantitativo de fenómenos económicos reales}. Su objetivo principal es dar \textbf{contenido empírico} a teorías cualitativas.

\subsection{Amalgama de Disciplinas}
\begin{enumerate}[leftmargin=8pt, noitemsep, topsep=2pt]
    \item \textbf{Teoría Económica:} Postula hipótesis cualitativas (e.g., relación precio-demanda).
    \item \textbf{Economía Matemática:} Expresa teorías en ecuaciones sin preocuparse por medición.
    \item \textbf{Estadística Económica:} Recolecta, procesa y presenta datos (PIB, empleo).
    \item \textbf{Estadística Matemática:} Proporciona herramientas de inferencia adaptadas a datos observacionales.
\end{enumerate}

\vogueQuote{La econometría transforma lo abstracto en mensurable, convirtiendo teoría en política.}

\section{Modelo de Regresión Lineal}
\subsection{Modelo Simple}
El modelo más elemental:

\vspace{2pt}
\eqbox{y_t = \beta_1 + \beta_2 X_t + u_t}
\vspace{2pt}

\begin{itemize}[leftmargin=8pt, noitemsep, topsep=2pt]
    \item $y_t$: Variable dependiente (regresando)
    \item $X_t$: Variable explicativa (regresor)
    \item $\beta_1, \beta_2$: Parámetros poblacionales
    \item $u_t$: \critical{Término de error estocástico}
\end{itemize}

\subsection{Modelo Matricial}
Representación compacta con $k$ regresores:

\vspace{2pt}
\eqbox{\mathbf{y} = \mathbf{X}\boldsymbol{\beta} + \mathbf{u}}
\vspace{2pt}

Donde $\mathbf{y}$ es $(n \times 1)$, $\mathbf{X}$ es $(n \times k)$, $\boldsymbol{\beta}$ es $(k \times 1)$, y $\mathbf{u}$ es $(n \times 1)$.

\subsection{El Término de Error}
El componente $u_t$ modela la \textbf{ignorancia parcial}:
\begin{itemize}[leftmargin=8pt, noitemsep, topsep=2pt]
    \item Variables omitidas
    \item Errores de medición
    \item Forma funcional incorrecta
    \item Aleatoriedad inherente
\end{itemize}

\critical{Supuesto clave:} $E(u_t|X_t) = 0$ → factores omitidos se cancelan en promedio.



\section{Supuestos Clásicos}
\subsection{Modelo Lineal Clásico}
\begin{tcolorbox}[colback=TitaniumGray, colframe=VogueBlack, arc=0pt, boxrule=0.3pt, title=\techfont \tiny SUPUESTOS GAUSS-MARKOV, fonttitle=\tiny\bfseries]
\scriptsize
\textbf{GM.1} Linealidad en parámetros \\
\textbf{GM.2} $E(u_t|X_t) = 0$ (media condicional cero) \\
\textbf{GM.3} $\text{Var}(u_t|X_t) = \sigma^2$ (homoscedasticidad) \\
\textbf{GM.4} $\text{Cov}(u_i, u_j|X) = 0, i \neq j$ (no autocorrelación) \\
\textbf{GM.5} $\text{Cov}(X_t, u_t) = 0$ (exogeneidad estricta) \\
\textbf{GM.6} $\text{rank}(\mathbf{X}) = k$ (no multicolinealidad perfecta)
\end{tcolorbox}

\subsection{Consecuencias}
Bajo GM.1-GM.6, OLS es \textbf{BLUE}: Best Linear Unbiased Estimator.

\section{Estimación: OLS}
\subsection{Principio de Mínimos Cuadrados}
Minimizar la suma de residuos al cuadrado:

\vspace{2pt}
$\min_{\boldsymbol{\beta}} \sum_{i=1}^{n}(y_i - \mathbf{x}_i'\boldsymbol{\beta})^2 = \min \mathbf{u}'\mathbf{u}$
\vspace{2pt}

\subsection{Estimador OLS}
\vspace{2pt}
\eqbox{\hat{\boldsymbol{\beta}} = (\mathbf{X}'\mathbf{X})^{-1}\mathbf{X}'\mathbf{y}}
\vspace{2pt}

\subsection{Propiedades en Muestra Finita}
\begin{enumerate}[leftmargin=8pt, noitemsep, topsep=2pt]
    \item \textbf{Insesgado:} $E(\hat{\boldsymbol{\beta}}) = \boldsymbol{\beta}$
    \item \textbf{Varianza:} $\text{Var}(\hat{\boldsymbol{\beta}}) = \sigma^2(\mathbf{X}'\mathbf{X})^{-1}$
    \item \textbf{Eficiente:} Mínima varianza entre estimadores lineales insesgados
\end{enumerate}

\subsection{Estimación de Varianza}
\vspace{2pt}
\eqbox{s^2 = \frac{\sum \hat{u}_i^2}{n-k} = \frac{\mathbf{u}'\mathbf{u}}{n-k}}
\vspace{2pt}

donde $\hat{u}_i = y_i - \mathbf{x}_i'\hat{\boldsymbol{\beta}}$ son los residuos.


% === ANIMACIÓN 1: OLS CONVERGENCIA (TAMAÑO EQUILIBRADO 2026) ===
\begin{center}
\vspace{6pt}
\subsection{Visualización OLS Dinámica}
\scriptsize\textit{Ajuste iterativo de la línea de regresión estructural}
\vspace{6pt}

\shorthandoff{>}\shorthandoff{<}
\begin{animateinline}[autoplay,loop]{8}
\multiframe{15}{rB=0.2+0.05}{
    % Scale 0.95 mantiene las fuentes legibles pero compactas
    \begin{tikzpicture}[scale=0.95]
    \begin{axis}[
        % Dimensiones equilibradas: 8cm x 5.5cm
        width=8cm, height=5.5cm,
        axis lines=middle,
        xlabel={\small $X$}, ylabel={\small $Y$},
        xmin=0, xmax=10, ymin=0, ymax=12,
        grid=major, grid style={dashed, gray!20},
        title={\small \techfont CONVERGENCIA OLS},
        title style={font=\small, yshift=3pt}
    ]
    \addplot[only marks, mark=*, mark size=1.4pt, NASAClassic]
        coordinates {(1,2) (2,3.5) (3,5) (4,6) (5,7.5) (6,9) (7,10) (8,11)};
    \addplot[Champagne, ultra thick, domain=0:10] {1 + \rB*x};
    \end{axis}
    \end{tikzpicture}
}
\end{animateinline}
\shorthandon{>}\shorthandon{<}
\vspace{6pt}
\end{center}



\section{Inferencia Estadística}
\subsection{Distribución del Estimador}
Bajo normalidad de errores ($u \sim N(0, \sigma^2\mathbf{I})$):
\eqbox{\hat{\boldsymbol{\beta}} \sim N\left(\boldsymbol{\beta}, \sigma^2(\mathbf{X}'\mathbf{X})^{-1}\right)}

\subsection{Test $t$ Individual}
Para probar $H_0: \beta_j = 0$ vs $H_1: \beta_j \neq 0$:
\eqbox{t = \frac{\hat{\beta}_j}{\text{se}(\hat{\beta}_j)} \sim t_{n-k}}

\critical{Regla de decisión:} Rechazar $H_0$ si $|t| > t_{\alpha/2, n-k}$.

\subsection{Test $F$ Conjunto}
Para probar $q$ restricciones lineales:
\eqbox{F = \frac{(\text{SSR}_r - \text{SSR}_{ur})/q}{\text{SSR}_{ur}/(n-k)} \sim F_{q, n-k}}

\subsection{Intervalos de Confianza}
Intervalo al $(1-\alpha)100\%$:
\eqbox{\hat{\beta}_j \pm t_{\alpha/2, n-k} \cdot \text{se}(\hat{\beta}_j)}

\section{Bondad de Ajuste}
\subsection{$R^2$ (Coeficiente de Determinación)}
\eqbox{R^2 = 1 - \frac{\text{SSR}}{\text{TSS}} = \frac{\text{ESS}}{\text{TSS}}}

\begin{itemize}[leftmargin=8pt, noitemsep, topsep=2pt]
    \item \textbf{TSS:} $\sum(y_i - \bar{y})^2$ (variación total)
    \item \textbf{ESS:} $\sum(\hat{y}_i - \bar{y})^2$ (variación explicada)
    \item \textbf{SSR:} $\sum\hat{u}_i^2$ (variación residual)
\end{itemize}

\critical{Limitación:} $R^2$ siempre aumenta con más regresores.

\subsection{$R^2$ Ajustado}
Penaliza la inclusión de variables irrelevantes que no aportan poder explicativo real al modelo:

\eqbox{\bar{R}^2 = 1 - \frac{\text{SSR}/(n-k)}{\text{TSS}/(n-1)}}


Puede \textit{decrecer} al agregar variables irrelevantes.

\section{Multicolinealidad}
\subsection{Problema}
Alta correlación entre regresores → $(\mathbf{X}'\mathbf{X})$ casi singular.

\subsection{Consecuencias}
\begin{itemize}[leftmargin=8pt, noitemsep, topsep=2pt]
    \item Varianzas infladas de $\hat{\boldsymbol{\beta}}$
    \item Estimadores inestables
    \item Tests $t$ individuales no significativos pero $F$ significativo
    \item Signos incorrectos
\end{itemize}

\subsection{Detección: VIF}
Factor de Inflación de Varianza:
\eqbox{\text{VIF}_j = \frac{1}{1 - R_j^2}}

donde $R_j^2$ es el $R^2$ de regresar $X_j$ sobre los demás regresores.

\critical{Regla práctica:} VIF $> 10$ indica multicolinealidad severa.

\subsection{Remedios}
\begin{enumerate}[leftmargin=8pt, noitemsep, topsep=2pt]
    \item Eliminar variables redundantes
    \item Aumentar tamaño muestral
    \item Recolectar nuevos datos
    \item Usar información \textit{a priori}
    \item Regresión Ridge (penalización)
\end{enumerate}

\section{Heterocedasticidad}
\subsection{Problema}
Varianza no constante: $\text{Var}(u_i|X_i) = \sigma_i^2$ (viola GM.3).

\subsection{Consecuencias}
\begin{itemize}[leftmargin=8pt, noitemsep, topsep=2pt]
    \item OLS sigue \textbf{insesgado} y \textbf{consistente}
    \item OLS \critical{NO es eficiente} (deja de ser BLUE)
    \item Errores estándar incorrectos
    \item Inferencia inválida (tests $t$ y $F$ no confiables)
\end{itemize}

\subsection{Detección}
\textbf{1. Test de White:}
\begin{itemize}[leftmargin=8pt, noitemsep, topsep=2pt]
    \item Regresión auxiliar: $\hat{u}_i^2 = \alpha_0 + \alpha_1 X_i + \alpha_2 X_i^2 + \alpha_3 X_i X_j + v_i$
    \item $H_0$: Homoscedasticidad
    \item Estadístico: $nR^2 \sim \chi^2_q$
\end{itemize}

\textbf{2. Test de Breusch-Pagan:}
\begin{itemize}[leftmargin=8pt, noitemsep, topsep=2pt]
    \item Regresión: $\hat{u}_i^2 = \gamma_0 + \gamma_1 Z_i + v_i$
    \item Similar a White pero con variables $Z$
\end{itemize}

\subsection{Soluciones}
\begin{enumerate}[leftmargin=8pt, noitemsep, topsep=2pt]
    \item \critical{Errores robustos} (White/Huber): Corrigen $\text{se}(\hat{\boldsymbol{\beta}})$
    \item \textbf{WLS:} Si se conoce forma de $\sigma_i^2$
    \item \textbf{FGLS:} Estimar $\sigma_i^2$ y aplicar WLS
\end{enumerate}

\textbf{Errores robustos de White:}
$$\widehat{\text{Var}}(\hat{\boldsymbol{\beta}}) = (\mathbf{X}'\mathbf{X})^{-1}\left(\sum_i \hat{u}_i^2\mathbf{x}_i\mathbf{x}_i'\right)(\mathbf{X}'\mathbf{X})^{-1}$$


% === ANIMACIÓN 2: HETEROCEDASTICIDAD (TAMAÑO REFINADO 2026) ===
\begin{center}
\vspace{4pt}
\subsection{Heterocedasticidad Visual}
\scriptsize\textit{Dispersión creciente de residuos}
\vspace{4pt}

\shorthandoff{>}\shorthandoff{<}
\begin{animateinline}[autoplay,loop]{6}
\multiframe{12}{rH=0.1+0.08}{
    % Scale 0.9 para un look más compacto y elegante
    \begin{tikzpicture}[scale=0.9]
    \begin{axis}[
        % Ancho reducido a 7.2cm para mayor equilibrio visual
        width=7.2cm, height=5cm,
        axis lines=middle,
        xlabel={\small $X$}, ylabel={\small $u$},
        xmin=0, xmax=10, ymin=-4, ymax=4,
        grid=major, grid style={dashed, gray!20},
        title={\small \techfont VAR$(u_i)$ CRECIENTE},
        title style={font=\small, yshift=2pt}
    ]
    \addplot[only marks, mark=*, mark size=1.3pt, CriticalRed]
        coordinates {
            (1,0.2) (2,-0.3) (3,0.5*\rH) (4,-0.7*\rH)
            (5,1.2*\rH) (6,-1.5*\rH) (7,2*\rH) (8,-2.3*\rH) (9,2.8*\rH)
        };
    \addplot[NASAClassic, dashed, domain=0:10, thick] {0};
    \end{axis}
    \end{tikzpicture}
}
\end{animateinline}
\shorthandon{>}\shorthandon{<}
\vspace{4pt}
\end{center}





\section{Autocorrelación}
\subsection{Problema}
Errores correlacionados en el tiempo: $\text{Cov}(u_i, u_j) \neq 0$ para $i \neq j$ (viola GM.4). Común en \textbf{series temporales}.

\subsection{Consecuencias}
Idénticas a heterocedasticidad:
\begin{itemize}[leftmargin=8pt, noitemsep, topsep=2pt]
    \item OLS insesgado pero no eficiente
    \item Errores estándar sesgados (generalmente subestimados)
    \item Intervalos de confianza y tests incorrectos
\end{itemize}

\subsection{Detección: Durbin-Watson}
\eqbox{DW = \frac{\sum_{t=2}^n(\hat{u}_t - \hat{u}_{t-1})^2}{\sum_{t=1}^n\hat{u}_t^2} \approx 2(1 - \hat{\rho})}

\begin{itemize}[leftmargin=8pt, noitemsep, topsep=2pt]
    \item DW $\approx 2$: No autocorrelación
    \item DW $< 2$: Autocorrelación positiva
    \item DW $> 2$: Autocorrelación negativa
\end{itemize}

\critical{Limitación:} Solo detecta autocorrelación de orden 1.

\subsection{Test de Breusch-Godfrey}
Más general que DW (permite múltiples rezagos y regresores estocásticos):
\begin{itemize}[leftmargin=8pt, noitemsep, topsep=2pt]
    \item Regresión auxiliar: $\hat{u}_t = \alpha_0 + \alpha_1 X_t + \rho_1 \hat{u}_{t-1} + \cdots + \rho_p \hat{u}_{t-p} + v_t$
    \item $(n-p)R^2 \sim \chi^2_p$ bajo $H_0$
\end{itemize}

\subsection{Soluciones}
\begin{enumerate}[leftmargin=8pt, noitemsep, topsep=2pt]
    \item \critical{Errores HAC} (Newey-West): Robustos a heterocedasticidad y autocorrelación
    \item \textbf{FGLS:} Estimar estructura de autocorrelación
    \item Modelar explícitamente: AR, MA, ARMA
\end{enumerate}

\section{Variables Instrumentales}
\subsection{Problema: Endogeneidad}
Si $\text{Cov}(X, u) \neq 0$, OLS es \critical{sesgado e inconsistente}.

\textbf{Causas:}
\begin{enumerate}[leftmargin=8pt, noitemsep, topsep=2pt]
    \item Variables omitidas correlacionadas con $X$
    \item Error de medición en $X$
    \item Simultaneidad (causalidad inversa)
    \item Variables dependientes rezagadas con errores AR
\end{enumerate}

\subsection{Solución: Instrumento $Z$}
Un instrumento válido satisface:
\begin{enumerate}[leftmargin=8pt, noitemsep, topsep=2pt]
    \item \textbf{Relevancia:} $\text{Cov}(Z, X) \neq 0$
    \item \textbf{Exogeneidad:} $\text{Cov}(Z, u) = 0$
\end{enumerate}

\subsection{Estimador IV (Simple)}
Con un regresor endógeno y un instrumento:
\eqbox{\hat{\beta}_{IV} = \frac{\text{Cov}(Z, Y)}{\text{Cov}(Z, X)} = \frac{\sum z_i y_i}{\sum z_i x_i}}

\subsection{Mínimos Cuadrados en Dos Etapas (2SLS)}
Para múltiples regresores/instrumentos:

\textbf{Etapa 1:} Regresión de $X$ sobre instrumentos $Z$:
$$X = Z\boldsymbol{\gamma} + v \quad \Rightarrow \quad \hat{X} = Z\hat{\boldsymbol{\gamma}}$$

\textbf{Etapa 2:} Regresión de $Y$ sobre $\hat{X}$:
$$Y = \hat{X}\boldsymbol{\beta} + u \quad \Rightarrow \quad \hat{\boldsymbol{\beta}}_{2SLS}$$

\subsection{Forma Matricial General}
\eqbox{\hat{\boldsymbol{\beta}}_{IV} = (\mathbf{X}'\mathbf{Z}(\mathbf{Z}'\mathbf{Z})^{-1}\mathbf{Z}'\mathbf{X})^{-1}\mathbf{X}'\mathbf{Z}(\mathbf{Z}'\mathbf{Z})^{-1}\mathbf{Z}'\mathbf{y}}

\subsection{Test de Hausman}
Para detectar endogeneidad:
$$H = (\hat{\boldsymbol{\beta}}_{OLS} - \hat{\boldsymbol{\beta}}_{IV})'\left[\widehat{\text{Var}}(\hat{\boldsymbol{\beta}}_{IV}) - \widehat{\text{Var}}(\hat{\boldsymbol{\beta}}_{OLS})\right]^{-1}(\hat{\boldsymbol{\beta}}_{OLS} - \hat{\boldsymbol{\beta}}_{IV})$$
$$H \sim \chi^2_k \text{ bajo } H_0: E(Xu) = 0$$

\vogueQuote{En economía, la causalidad es un lujo que solo los instrumentos pueden comprar.}

\section{Teoría Asintótica}
\subsection{Motivación}
Cuando fallan supuestos clásicos (regresores estocásticos, errores no normales), las propiedades de muestra finita se pierden. Recurrimos a \textbf{propiedades asintóticas} ($n \to \infty$).

\subsection{Conceptos de Convergencia}
\textbf{1. Convergencia en Probabilidad ($\xrightarrow{p}$):}
$$\hat{\theta}_n \xrightarrow{p} \theta \iff \lim_{n\to\infty} P(|\hat{\theta}_n - \theta| > \epsilon) = 0$$

\textbf{2. Convergencia Casi Segura ($\xrightarrow{a.s.}$):}
$$P\left(\lim_{n\to\infty}\hat{\theta}_n = \theta\right) = 1$$

\textbf{3. Convergencia en Distribución ($\xrightarrow{d}$):}
$$F_{\hat{\theta}_n}(x) \to F_\theta(x) \text{ en puntos de continuidad}$$

\subsection{Ley de Grandes Números (LLN)}
Si $X_1, \ldots, X_n$ i.i.d. con $E(X_i) = \mu < \infty$:
\eqbox{\bar{X}_n = \frac{1}{n}\sum_{i=1}^n X_i \xrightarrow{p} \mu}

\textbf{Versiones:}
\begin{itemize}[leftmargin=8pt, noitemsep, topsep=2pt]
    \item \textbf{Kolmogorov:} Para i.i.d. (convergencia a.s.)
    \item \textbf{Markov:} Para i.n.i.d. (heterogéneos)
    \item \textbf{Ergódico:} Para series dependientes
\end{itemize}

\critical{Implicación:} Promedios muestrales son estimadores consistentes.

\subsection{Teorema del Límite Central (CLT)}
Si $X_1, \ldots, X_n$ i.i.d. con $E(X_i) = \mu$, $\text{Var}(X_i) = \sigma^2 < \infty$:
\eqbox{\sqrt{n}(\bar{X}_n - \mu) \xrightarrow{d} N(0, \sigma^2)}

\textbf{Versiones Generales:}
\begin{itemize}[leftmargin=8pt, noitemsep, topsep=2pt]
    \item \textbf{Lindeberg-Levy:} Para i.i.d.
    \item \textbf{Lindeberg-Feller:} Para i.n.i.d.
    \item \textbf{FCLT:} Para procesos dependientes (movimiento browniano)
\end{itemize}

\subsection{Propiedades Asintóticas de OLS}
\textbf{Consistencia:}
$\hat{\boldsymbol{\beta}} \xrightarrow{p} \boldsymbol{\beta}$
requiere: $\frac{1}{n}\mathbf{X}'\mathbf{X} \xrightarrow{p} Q$ (definida positiva) y $\frac{1}{n}\mathbf{X}'\mathbf{u} \xrightarrow{p} 0$.

\textbf{Normalidad Asintótica:}\\[10pt]
\eqbox{\sqrt{n}(\hat{\boldsymbol{\beta}} - \boldsymbol{\beta}) \xrightarrow{d} N(0, \sigma^2 Q^{-1})}

donde $Q = \text{plim}\frac{1}{n}\mathbf{X}'\mathbf{X}$.

% === ANIMACIÓN 3: CONVERGENCIA ASINTÓTICA (TAMAÑO REFINADO 2026) ===
\begin{center}
\vspace{4pt}
\subsection{Convergencia del Estimador}
\scriptsize\textit{Distribución de $\hat{\beta}$ concentrándose a medida que aumenta $n$}
\vspace{6pt}

\shorthandoff{>}\shorthandoff{<}
\begin{animateinline}[autoplay,loop]{5}
\multiframe{10}{rN=50+50}{
    % Scale 0.9 para consistencia con el gráfico anterior
    \begin{tikzpicture}[scale=0.9]
    \begin{axis}[
        % Dimensiones espejo: 7.2cm x 5cm
        width=7.2cm, height=5cm,
        axis lines=middle,
        xlabel={\small $\hat{\beta}$}, ylabel={\small Densidad},
        xmin=-1, xmax=3, ymin=0, ymax=2.5,
        grid=major, grid style={dashed, gray!20},
        title={\small \techfont MUESTRA $n = \rN$},
        title style={font=\small, yshift=2pt}
    ]
    \addplot[NASAClassic, thick, domain=-1:3, samples=50]
        {exp(-((x-1)^2)/(2*(1/sqrt(\rN/50))))/sqrt(2*pi*(1/sqrt(\rN/50)))};
    \addplot[CriticalRed, dashed, thick] coordinates {(1,0) (1,2.5)};
    \end{axis}
    \end{tikzpicture}
}
\end{animateinline}
\shorthandon{>}\shorthandon{<}
\vspace{4pt}
\end{center}


\section{GMM: Método Generalizado de Momentos}
\subsection{Principio}
Estimar parámetros igualando momentos poblacionales a muestrales:
$E[g(W_i, \boldsymbol{\theta}_0)] = 0 \quad (\text{condiciones de momento})$

\subsection{Momento Muestral}
$\bar{g}_n(\boldsymbol{\theta}) = \frac{1}{n}\sum_{i=1}^n g(W_i, \boldsymbol{\theta})$

\subsection{Estimador GMM}
Minimizar forma cuadrática:
\eqbox{\hat{\boldsymbol{\theta}}_{GMM} = \arg\min_{\boldsymbol{\theta}} \bar{g}_n(\boldsymbol{\theta})'\mathbf{W}_n\bar{g}_n(\boldsymbol{\theta})}

donde $\mathbf{W}_n$ es matriz de ponderación.

\subsection{GMM Óptimo}
\critical{Matriz óptima:} $\mathbf{W}_n = \hat{\mathbf{V}}_n^{-1}$ donde
$\mathbf{V}_n = \text{Var}(\sqrt{n}\bar{g}_n) \xrightarrow{p} \mathbf{V}$

\textbf{Resultado:} Estimador asintóticamente eficiente.

\subsection{Test de Sobreidentificación (J)}
Si $\#\text{momentos} > \#\text{parámetros}$:
$J = n\bar{g}_n(\hat{\boldsymbol{\theta}})'\mathbf{W}_n\bar{g}_n(\hat{\boldsymbol{\theta}}) \sim \chi^2_{r-k}$



\section{Series de Tiempo}
\subsection{Proceso Estocástico}
Secuencia $\{Y_t: t \in T\}$ de variables aleatorias.

\textbf{Realización:} Valores observados $\{y_1, \ldots, y_T\}$.

\subsection{Estacionaridad}
\textbf{Fuerte:} Distribución invariante a desplazamientos temporales.

\textbf{Débil (Covarianza):}
\begin{enumerate}[leftmargin=8pt, noitemsep, topsep=2pt]
    \item $E(Y_t) = \mu$ (constante)
    \item $\text{Var}(Y_t) = \sigma^2$ (constante)
    \item $\text{Cov}(Y_t, Y_{t-s}) = \gamma_s$ (depende solo de $s$)
\end{enumerate}

\subsection{Proceso AR(1)}
\eqbox{Y_t = \phi Y_{t-1} + \epsilon_t, \quad \epsilon_t \sim WN(0, \sigma^2)}

\textbf{Estacionario si:} $|\phi| < 1$

\textbf{Momentos:}
\begin{itemize}[leftmargin=8pt, noitemsep, topsep=2pt]
    \item $E(Y_t) = 0$ (si $\epsilon_t$ tiene media cero)
    \item $\text{Var}(Y_t) = \frac{\sigma^2}{1-\phi^2}$
    \item $\text{Cov}(Y_t, Y_{t-s}) = \phi^s \frac{\sigma^2}{1-\phi^2}$
\end{itemize}

\subsection{Raíz Unitaria}
Si $\phi = 1$: \critical{Caminata aleatoria} (no estacionaria).
$Y_t = Y_{t-1} + \epsilon_t \Rightarrow Y_t = Y_0 + \sum_{i=1}^t \epsilon_i$

\textbf{Propiedades:}
\begin{itemize}[leftmargin=8pt, noitemsep, topsep=2pt]
    \item $\text{Var}(Y_t) = t\sigma^2$ (crece con $t$)
    \item Shocks tienen efecto permanente
    \item Proceso $I(1)$ (integrado de orden 1)
\end{itemize}

\subsection{Test de Dickey-Fuller}
Probar $H_0$: raíz unitaria
$\Delta Y_t = \delta Y_{t-1} + \epsilon_t$

donde $\delta = \phi - 1$. $H_0: \delta = 0$.

\critical{Distribución no estándar:} Usa valores críticos de Fuller.

\textbf{ADF (Aumentado):} Agrega rezagos de $\Delta Y_t$ para capturar autocorrelación.

\subsection{Cointegración}
Dos series $I(1)$, $X_t$ e $Y_t$, están cointegradas si:
\eqbox{Y_t - \beta X_t \sim I(0)}

\critical{Combinación lineal es estacionaria.}

\textbf{Test de Engle-Granger:}
\begin{enumerate}[leftmargin=8pt, noitemsep, topsep=2pt]
    \item Estimar $Y_t = \alpha + \beta X_t + u_t$ por OLS
    \item Test ADF sobre residuos $\hat{u}_t$
    \item Si $\hat{u}_t \sim I(0)$, las series están cointegradas
\end{enumerate}

% === ANIMACIÓN 4: PROCESO AR(1) (TAMAÑO REFINADO 2026) ===
\begin{center}
\vspace{4pt}
\subsection{Proceso AR(1) Dinámico}
\scriptsize\textit{Evolución de una serie temporal estacionaria en tiempo real}
\vspace{6pt}

\shorthandoff{>}\shorthandoff{<}
\begin{animateinline}[autoplay,loop]{6}
\multiframe{20}{rT=1+1}{
    % Scale 0.9 para consistencia editorial 2026
    \begin{tikzpicture}[scale=0.9]
    \begin{axis}[
        % Tamaño optimizado: 7.2cm x 5cm
        width=7.2cm, height=5cm,
        axis lines=middle,
        xlabel={\small $t$}, ylabel={\small $Y_t$},
        xmin=0, xmax=20, ymin=-3, ymax=3,
        grid=major, grid style={dashed, gray!20},
        title={\small \techfont DINÁMICA AR(1) CON $\phi=0.7$},
        title style={font=\small, yshift=2pt}
    ]
    \addplot[NASAClassic, thick, mark=*, mark size=1pt]
        coordinates {
            (1,0) (2,0.5) (3,0.8) (4,0.6) (5,0.9)
            (6,1.1) (7,0.7) (8,0.4) (9,0.6) (10,0.3)
            (11,-0.1) (12,0.2) (13,0.5) (14,0.8) (15,0.6)
            (16,0.4) (17,0.7) (18,0.9) (19,0.5) (20,0.3)
        };
    \addplot[CriticalRed, dashed, thick] coordinates {(\rT,-3) (\rT,3)};
    \end{axis}
    \end{tikzpicture}
}
\end{animateinline}
\shorthandon{>}\shorthandon{<}
\vspace{4pt}
\end{center}


\section{Tipos de Datos}
\begin{tcolorbox}[colback=TitaniumGray, colframe=NASAClassic, arc=0pt, boxrule=0.3pt, title=\techfont \tiny SERIES DE TIEMPO, fonttitle=\tiny\bfseries]
\scriptsize
Observaciones a lo largo del tiempo: $Y_t$ para $t=1,\ldots,T$.

\critical{Problemas:} Autocorrelación, no estacionaridad, raíz unitaria.

\textbf{Ejemplo:} PIB trimestral 2000-2025.
\end{tcolorbox}

\begin{tcolorbox}[colback=TitaniumGray, colframe=NASAClassic, arc=0pt, boxrule=0.3pt, title=\techfont \tiny CORTE TRANSVERSAL, fonttitle=\tiny\bfseries]
\scriptsize
Observaciones de múltiples unidades en $t$ fijo: $Y_i$ para $i=1,\ldots,N$.

\critical{Problemas:} Heterogeneidad, heterocedasticidad.

\textbf{Ejemplo:} Ingresos de 1000 hogares en 2026.
\end{tcolorbox}

\begin{tcolorbox}[colback=TitaniumGray, colframe=NASAClassic, arc=0pt, boxrule=0.3pt, title=\techfont \tiny DATOS DE PANEL, fonttitle=\tiny\bfseries]
\scriptsize
Combinación: $Y_{it}$ para $i=1,\ldots,N$ y $t=1,\ldots,T$.

\critical{Ventaja:} Controla heterogeneidad no observada.

\textbf{Ejemplo:} 500 hogares durante 10 años.
\end{tcolorbox}



\section{Modelos de Panel}
\subsection{Efectos Fijos (FE)}
\eqbox{Y_{it} = \alpha_i + \mathbf{X}_{it}'\boldsymbol{\beta} + u_{it}}

$\alpha_i$: Efecto individual fijo (heterogeneidad no observada).

\textbf{Estimación:} Transformación \textit{within} (demeaning).

\subsection{Efectos Aleatorios (RE)}
\eqbox{Y_{it} = \alpha + \mathbf{X}_{it}'\boldsymbol{\beta} + (\eta_i + u_{it})}

$\eta_i$: Efecto aleatorio (i.i.d.).

\textbf{Supuesto clave:} $E(\eta_i|\mathbf{X}_{it}) = 0$.

\subsection{Test de Hausman}
Decidir entre FE y RE:
$H = (\hat{\boldsymbol{\beta}}_{FE} - \hat{\boldsymbol{\beta}}_{RE})'\left[\widehat{\text{Var}}(\hat{\boldsymbol{\beta}}_{FE}) - \widehat{\text{Var}}(\hat{\boldsymbol{\beta}}_{RE})\right]^{-1}(\hat{\boldsymbol{\beta}}_{FE} - \hat{\boldsymbol{\beta}}_{RE})$

$H_0$: RE es consistente (y eficiente). Si se rechaza, usar FE.

\section{Modelos No Lineales}
\subsection{Modelo Probit}
Para variable dependiente binaria $Y \in \{0,1\}$:
\eqbox{P(Y=1|\mathbf{X}) = \Phi(\mathbf{X}'\boldsymbol{\beta})}

$\Phi$: CDF normal estándar.

\subsection{Modelo Logit}
\eqbox{P(Y=1|\mathbf{X}) = \frac{\exp(\mathbf{X}'\boldsymbol{\beta})}{1 + \exp(\mathbf{X}'\boldsymbol{\beta})} = \Lambda(\mathbf{X}'\boldsymbol{\beta})}

$\Lambda$: Función logística.

\subsection{Modelo Tobit}
Para datos censurados:
$Y_i^* = \mathbf{X}_i'\boldsymbol{\beta} + u_i$
$Y_i = \begin{cases} Y_i^* & \text{si } Y_i^* > 0 \\ 0 & \text{si } Y_i^* \leq 0 \end{cases}$

\textbf{Estimación:} Máxima verosimilitud.

\section{Máxima Verosimilitud}
\subsection{Principio}
Maximizar función de verosimilitud:
$L(\boldsymbol{\theta}; \mathbf{y}) = \prod_{i=1}^n f(y_i|\boldsymbol{\theta})$

O log-verosimilitud:
\eqbox{\ell(\boldsymbol{\theta}) = \sum_{i=1}^n \ln f(y_i|\boldsymbol{\theta})}

\subsection{Estimador ML}
\eqbox{\hat{\boldsymbol{\theta}}_{ML} = \arg\max_{\boldsymbol{\theta}} \ell(\boldsymbol{\theta})}

\textbf{CPO:} $\frac{\partial \ell}{\partial \boldsymbol{\theta}} = 0$

\subsection{Propiedades Asintóticas}
\begin{enumerate}[leftmargin=8pt, noitemsep, topsep=2pt]
    \item \textbf{Consistencia:} $\hat{\boldsymbol{\theta}}_{ML} \xrightarrow{p} \boldsymbol{\theta}_0$
    \item \textbf{Normalidad asintótica:}
    $\sqrt{n}(\hat{\boldsymbol{\theta}}_{ML} - \boldsymbol{\theta}_0) \xrightarrow{d} N(0, \mathcal{I}^{-1})$
    \item \textbf{Eficiencia:} Alcanza cota Cramér-Rao
\end{enumerate}

\textbf{Matriz de Información:}
$\mathcal{I}(\boldsymbol{\theta}) = -E\left[\frac{\partial^2 \ell}{\partial \boldsymbol{\theta}\partial\boldsymbol{\theta}'}\right]$

\section{Tests Avanzados}
\subsection{Test de Wald}
Para $H_0: \mathbf{R}\boldsymbol{\beta} = \mathbf{r}$:
\eqbox{W = (\mathbf{R}\hat{\boldsymbol{\beta}} - \mathbf{r})'\left[\mathbf{R}\widehat{\text{Var}}(\hat{\boldsymbol{\beta}})\mathbf{R}'\right]^{-1}(\mathbf{R}\hat{\boldsymbol{\beta}} - \mathbf{r})}
$W \sim \chi^2_q$

\subsection{Test de Razón de Verosimilitud (LR)}
\eqbox{LR = 2[\ell(\hat{\boldsymbol{\theta}}_{ur}) - \ell(\hat{\boldsymbol{\theta}}_r)] \sim \chi^2_q}

\subsection{Test LM (Multiplicador de Lagrange)}
Basado en gradiente en el restringido:
$LM = \mathbf{s}(\hat{\boldsymbol{\theta}}_r)'\mathcal{I}^{-1}(\hat{\boldsymbol{\theta}}_r)\mathbf{s}(\hat{\boldsymbol{\theta}}_r) \sim \chi^2_q$

\critical{Equivalencia:} Wald, LR y LM son asintóticamente equivalentes.

\section{Metodología Clásica}
\subsection{8 Pasos (Gujarati)}
\begin{enumerate}[leftmargin=8pt, itemsep=1pt, topsep=2pt]
    \item \textbf{Teoría/Hipótesis}
    \item \textbf{Modelo Matemático}
    \item \textbf{Modelo Econométrico} (agregar $u_t$)
    \item \textbf{Obtención de Datos}
    \item \textbf{Estimación de Parámetros}
    \item \textbf{Prueba de Hipótesis}
    \item \textbf{Pronóstico}
    \item \textbf{Política Económica}
\end{enumerate}

\section{Batería de Tests}
\vspace{4pt}
\begin{center}
\scriptsize\techfont
\begin{tabular}{@{}lll@{}}
\toprule
\textbf{Problema} & \textbf{Test} & \textbf{$H_0$} \\
\midrule
Multicolinealidad & VIF & VIF $< 10$ \\
Heterocedasticidad & White & Homosced. \\
 & Breusch-Pagan & Homosced. \\
Autocorrelación & Durbin-Watson & $\rho = 0$ \\
 & Breusch-Godfrey & No autocorr. \\
Normalidad & Jarque-Bera & Normal \\
Especificación & RESET & Correcta \\
Raíz Unitaria & ADF & Raíz unit. \\
 & PP & Raíz unit. \\
Cointegración & Engle-Granger & No cointegr. \\
Endogeneidad & Hausman & Exogeneidad \\
\bottomrule
\end{tabular}
\end{center}
\vspace{4pt}

\section{Fórmulas Esenciales}
\begin{tcolorbox}[colback=NASAClassic!10, colframe=NASAClassic, arc=0pt, boxrule=0.4pt, title=\techfont \tiny OLS, fonttitle=\tiny\bfseries]
\scriptsize
$\hat{\boldsymbol{\beta}} = (\mathbf{X}'\mathbf{X})^{-1}\mathbf{X}'\mathbf{y}$
$\text{Var}(\hat{\boldsymbol{\beta}}) = \sigma^2(\mathbf{X}'\mathbf{X})^{-1}$
$s^2 = \frac{\mathbf{u}'\mathbf{u}}{n-k}$
\end{tcolorbox}

\begin{tcolorbox}[colback=Champagne!15, colframe=Champagne, arc=0pt, boxrule=0.4pt, title=\techfont \tiny IV/2SLS, fonttitle=\tiny\bfseries]
\scriptsize
$\hat{\boldsymbol{\beta}}_{IV} = (\mathbf{X}'\mathbf{Z}(\mathbf{Z}'\mathbf{Z})^{-1}\mathbf{Z}'\mathbf{X})^{-1}\mathbf{X}'\mathbf{Z}(\mathbf{Z}'\mathbf{Z})^{-1}\mathbf{Z}'\mathbf{y}$
\end{tcolorbox}

\begin{tcolorbox}[colback=CriticalRed!10, colframe=CriticalRed, arc=0pt, boxrule=0.4pt, title=\techfont \tiny GMM, fonttitle=\tiny\bfseries]
\scriptsize
$\hat{\boldsymbol{\theta}} = \arg\min \bar{g}_n'\mathbf{W}_n\bar{g}_n$
$\mathbf{W}_n^* = \hat{\mathbf{V}}_n^{-1}$
\end{tcolorbox}

\vogueQuote{La sofisticación estadística sin intuición económica es mera pirotecnia matemática.}

\section{Guía Práctica}
\subsection{Checklist de Regresión}
\begin{enumerate}[leftmargin=8pt, itemsep=1pt, topsep=2pt]
    \item[$\square$] Graficar datos
    \item[$\square$] Verificar estacionaridad (series)
    \item[$\square$] Estimar modelo OLS
    \item[$\square$] Analizar residuos
    \item[$\square$] Tests diagnósticos
    \item[$\square$] Corregir violaciones
    \item[$\square$] Validar externamente
    \item[$\square$] Interpretar económicamente
\end{enumerate}

\subsection{Errores Comunes}
\begin{itemize}[leftmargin=8pt, noitemsep, topsep=2pt]
    \item \critical{NO} comparar $R^2$ con diferentes $Y$
    \item \critical{NO} confundir correlación con causalidad
    \item \critical{NO} ignorar errores robustos
    \item \critical{NO} usar OLS con endogeneidad
    \item \critical{NO} olvidar supuestos
\end{itemize}

\subsection{Software}
\begin{itemize}[leftmargin=8pt, noitemsep, topsep=2pt]
    \item \textbf{R:} \texttt{lm()}, \texttt{plm}, \texttt{tseries}
    \item \textbf{Python:} \texttt{statsmodels}, \texttt{linearmodels}
    \item \textbf{Stata:} \texttt{reg}, \texttt{ivregress}, \texttt{xtset}
\end{itemize}

\section{Referencias Core}
\begin{itemize}[leftmargin=8pt, noitemsep, topsep=2pt]
    \item \textbf{Gujarati \& Porter} (2009). \textit{Econometría}. McGraw-Hill. 5ed.
    \item \textbf{Wooldridge} (2015). \textit{Introductory Econometrics}. Cengage. 6ed.
    \item \textbf{Davidson \& MacKinnon} (2004). \textit{Econometric Theory and Methods}. Oxford.
    \item \textbf{Hamilton} (1994). \textit{Time Series Analysis}. Princeton.
    \item \textbf{Greene} (2018). \textit{Econometric Analysis}. Pearson. 8ed.
\end{itemize}

\end{multicols*}

% === PIE DE PÁGINA FINAL ===
\vfill
\hrule height 0.3pt
\vspace{3pt}
\begin{center}
    \scriptsize\techfont\lsstyle
    \textcolor{VogueBlack}{ECONOMETRÍA 2026} \textbar\
    \textcolor{Champagne}{EMANUEL QUINTANA SILVA} \textbar\
    UPTC ECONOMICS \\
    \vspace{2pt}
    \textcolor{InkGray}{\href{mailto:emanuel.quintana@uptc.edu.co}{emanuel.quintana@uptc.edu.co} \textbar\ ORCID: 0009-0006-8419-2805 \textbar\ GitHub: @emanuelquintana-glitch} \\
    \vspace{2pt}
    \textit{Compilado con pdfLaTeX + animate. Abrir con Adobe Acrobat Reader para animaciones o para usuarios linux firefox.}
\end{center}

\end{document}
